\documentclass[10pt,a4paper]{article}
\usepackage[utf8]{inputenc}
\usepackage[vietnamese]{babel}
\usepackage{amsmath,amssymb,amsfonts}
\usepackage{geometry}
\geometry{a4paper, left=10mm, right=10mm, top=10mm, bottom=10mm}
\usepackage{tikz}
\usetikzlibrary{shapes,arrows.meta,positioning,calc}
\usepackage[most]{tcolorbox}
\usepackage{array,booktabs,multicol}
\usepackage{fontawesome}
\usepackage{helvet}
\renewcommand{\familydefault}{\sfdefault}
\usepackage{parskip}

\definecolor{primary}{RGB}{112, 48, 160}
\definecolor{teal}{RGB}{0, 128, 128}
\definecolor{orange}{RGB}{230, 108, 44}
\definecolor{green}{RGB}{46, 139, 87}
\definecolor{rose}{RGB}{192, 0, 0}

\pagestyle{empty}

\begin{document}

% --- HEADER BANNER ---
\begin{tcolorbox}[
  enhanced,
  colback=violet!8!white,
  colframe=violet!80!black,
  arc=2.5mm, boxrule=1pt,
  left=3mm, right=3mm, top=2mm, bottom=2mm
]
  \begin{tabular*}{\textwidth}{@{\extracolsep{\fill}} lr}
    {\small\bfseries\color{violet!85!black}\faBook\ SỔ TAY KIẾN THỨC TOÁN 2} & 
    {\small\bfseries\color{teal!85!black}\faGraduationCap\ CHƯƠNG TRÌNH GDPT 2018}
  \end{tabular*}
  \vspace{-1mm}
  \begin{center}
    \rule{0.98\textwidth}{0.6pt}
  \end{center}
  \vspace{-2mm}
  \centering
  {\LARGE\bfseries\color{violet!90!black}\faBookmark\ BÀI 57: THỰC HÀNH VÀ TRẢI NGHIỆM ĐO ĐỘ DÀI}\\[1mm]
  {\itshape\footnotesize\color{black!75} Kĩ năng sử dụng thước đo, ước lượng và thực hành đo đạc đồ vật xung quanh}
\end{tcolorbox}

\vspace{-2mm}

\noindent
\begin{minipage}[t]{0.49\textwidth}\vspace{0pt}
\begin{tcolorbox}[
  enhanced,
  colback=violet!4!white,
  colframe=violet!75!black,
  arc=2mm, boxrule=0.8pt,
  left=2mm, right=2mm, top=1.8mm, bottom=1.8mm,
  title={\footnotesize\bfseries\color{white}\faKey\ 1. Dụng cụ đo và cách sử dụng thước}
]
{\footnotesize
\begin{itemize}\setlength{\itemsep}{1pt}\setlength{\leftskip}{-2.5mm}
  \item \textbf{Các dụng cụ đo thường dùng:}\\
  Thước kẻ học sinh, thước dây, thước cuộn, thước gấp.
  \item \textbf{Các bước đo độ dài đồ vật bằng thước kẻ:}
  \begin{itemize}\setlength{\itemsep}{0.5pt}\setlength{\leftskip}{-2mm}
    \item \textbf{Bước 1:} Đặt vạch số \textbf{0} của thước trùng với một đầu của vật cần đo.
    \item \textbf{Bước 2:} Giữ thước áp sát và thẳng theo mép vật.
    \item \textbf{Bước 3:} Đầu kia của vật trùng với vạch số nào thì đọc và ghi số đo kèm theo đơn vị.
  \end{itemize}
\end{itemize}
}
\end{tcolorbox}
\end{minipage}
\hfill
\begin{minipage}[t]{0.49\textwidth}\vspace{0pt}
\begin{tcolorbox}[
  enhanced,
  colback=teal!4!white,
  colframe=teal!75!black,
  arc=2mm, boxrule=0.8pt,
  left=2mm, right=2mm, top=1.8mm, bottom=1.8mm,
  title={\footnotesize\bfseries\color{white}\faSliders\ 2. Kĩ năng ước lượng độ dài}
]
{\footnotesize
\begin{itemize}\setlength{\itemsep}{1pt}\setlength{\leftskip}{-2.5mm}
  \item \textbf{Ước lượng là gì?} Là dùng mắt và kinh nghiệm để đoán độ dài của một vật trước khi đo thực tế.
  \item \textbf{Các mốc quy chiếu quen thuộc để ước lượng:}
  \begin{itemize}\setlength{\itemsep}{0.5pt}\setlength{\leftskip}{-2mm}
    \item Độ rộng ngón tay trỏ: khoảng $1\ cm$.
    \item Một gang tay học sinh: khoảng $13 - 15\ cm$ (khoảng $1{,}5\ dm$).
    \item Một bước chân học sinh: khoảng $40 - 50\ cm$.
    \item Sải tay mở rộng của học sinh: khoảng $1\ m$.
  \end{itemize}
  \item So sánh giữa số đo ước lượng và số đo thực tế để rèn luyện độ chính xác.
\end{itemize}
}
\end{tcolorbox}
\end{minipage}

\vspace{-1.5mm}
\begin{tcolorbox}[
  enhanced,
  colback=orange!4!white,
  colframe=orange!80!black,
  arc=2mm, boxrule=0.8pt,
  left=2mm, right=2mm, top=1.5mm, bottom=1.5mm,
  title={\footnotesize\bfseries\color{white}\faSitemap\ 3. Quy trình 3 bước thực hành đo độ dài chuẩn mực}
]
\centering
\begin{tikzpicture}[x=1.65cm, y=0.75cm]
  \tikzset{arr/.style={-{Stealth[scale=0.8]}}, every node/.append style={align=center}}
  \node[rectangle, rounded corners=2pt, draw=violet!80!black, fill=violet!15, font=\scriptsize\bfseries] (b1) at (0,0) {BƯỚC 1: ƯỚC LƯỢNG\\Nhìn và đoán số đo};
  \node[rectangle, rounded corners=2pt, draw=teal!80!black, fill=teal!15, font=\scriptsize\bfseries] (b2) at (2.4,0) {BƯỚC 2: ĐẶT THƯỚC\\Đầu vật trùng vạch 0};
  \node[rectangle, rounded corners=2pt, draw=orange!80!black, fill=orange!15, font=\scriptsize\bfseries] (b3) at (4.8,0) {BƯỚC 3: ĐỌC \& GHI\\Kèm theo tên đơn vị};

  \draw[arr, thick] (b1) -- (b2);
  \draw[arr, thick] (b2) -- (b3);
\end{tikzpicture}
\end{tcolorbox}

\vspace{-1.5mm}
\noindent
\begin{minipage}[t]{0.49\textwidth}\vspace{0pt}
\begin{tcolorbox}[
  enhanced,
  colback=green!4!white,
  colframe=green!60!black,
  arc=2mm, boxrule=0.8pt,
  left=2mm, right=2mm, top=1.8mm, bottom=1.8mm,
  title={\footnotesize\bfseries\color{white}\faCheckCircle\ 4. Bảng ghi chép thực hành mẫu}
]
{\footnotesize
\begin{tabular}{|l|c|c|}
  \hline
  \textbf{Tên đồ vật} & \textbf{Ước lượng} & \textbf{Đo thực tế} \\
  \hline
  Bút chì & $15\ cm$ & $16\ cm$ \\
  Hộp bút & $2\ dm$ & $21\ cm$ \\
  Bề dài bàn học & $1\ m$ & $110\ cm$ \\
  Chiều dài lớp học & $7\ m$ & $7\ m\ 20\ cm$ \\
  \hline
\end{tabular}
\vspace{1mm}
\textit{Nhận xét: Số đo ước lượng gần sát với kết quả đo thực tế.}
}
\end{tcolorbox}
\end{minipage}
\hfill
\begin{minipage}[t]{0.49\textwidth}\vspace{0pt}
\begin{tcolorbox}[
  enhanced,
  colback=rose!4!white,
  colframe=rose!70!black,
  arc=2mm, boxrule=0.8pt,
  left=2mm, right=2mm, top=1.8mm, bottom=1.8mm,
  title={\footnotesize\bfseries\color{white}\faExclamationTriangle\ 5. Lỗi đo lường thường gặp}
]
{\footnotesize
\begin{itemize}\setlength{\itemsep}{1.5pt}\setlength{\leftskip}{-2.5mm}
  \item \textbf{Đặt đầu vật vào mép thước:} Rất nhiều thước kẻ có khoảng cách từ mép thước đến vạch số $0$. Đặt đầu vật vào mép ngoài là bị \textbf{sai số}.
  \item \textbf{Đặt thước bị lệch, chéo:} Thước phải đặt song song và áp sát mép của vật.
  \item Mắt nhìn thẳng góc với vạch chia của thước để đọc số đo chuẩn nhất.
\end{itemize}
}
\end{tcolorbox}
\end{minipage}


\end{document}
